\documentclass[12pt,a4paper]{article}
\usepackage{ucebnice}
\usepackage{polynom}

\title{Exponenciální rovnice a jejich soustavy}
\author{Ondřej Škubala}
\date{\today}

\begin{document}

\maketitle
\newpage

\section{Exponenciální rovnice}

\subsection{Exponenciální rovnice}

Exponenciální rovnice jsou takové rovnice, v níž se neznámá nevyskytuje
v základní části výrazu, ale v jeho \important{exponentu}. Tato skutečnost výrazně
ovlivňuje nejen způsob, jakým rovnici upravujeme, ale i metody, které při jejím řešení používáme.

\medskip

Obecně můžeme exponenciální rovnici zapsat ve tvaru
\[
    a^{f(x)} = b^{g(x)},
\]
kde $a>0$, $a\neq1$ a také $b>0$. Výrazy $f(x)$ a $g(x)$ bývají nejčastěji lineární
nebo kvadratické funkce.

\begin{examplebox}
Uvažujme výraz
\[
    2^x = 8.
\]
Protože $8 = 2^3$, je zřejmé, že řešením je právě $x=3$.  
Tedy kořenem této rovnice je množina $K=\{3\}$.
\end{examplebox}

Tento jednoduchý příklad jasně ukazuje hlavní princip exponenciálních rovnic:

\begin{notebox}
Hledáme takovou hodnotu exponentu, která při umocnění daného základu 
vede k požadovanému výsledku.
\end{notebox}

Tuto situaci si lze velmi dobře představit i graficky. Jak již víme, tak $y = 2^x$ je exponenciální
křivka a rovnice $2^x = 8$ tedy znamená hledat \textbf{průsečík} grafu exponenciální funkce
s lineární přímkou $y = 8$.

\begin{notebox}
Řešení exponenciální rovnice lze chápat jako hledání bodu, 
ve kterém exponenciální křivka protne danou vodorovnou přímku.
\end{notebox}

\begin{center}
\begin{axisgraph}[
    xmin=-2, xmax=5,
    ymin=0, ymax=10,
    width=10cm,
    height=6cm,
    samples=200,
    domain=-2:5,
    title={Grafické znázornění rovnice $2^x = 8$}
]
    % exponenciála
    \addplot[blue, thick] {2^x};

    % přímka y=8
    \addplot[red, dashed, thick] {8};

    % průsečík (3,8)
    \addplot[
        mark=*,
        mark options={scale=1.2, fill=black},
    ] coordinates {(3,8)};

\end{axisgraph}
\end{center}

\subsubsection*{Klíčová vlastnost exponenciální funkce}

Při řešení exponenciálních rovnic využíváme toho, že funkce $y=a^x$ je pro každý
kladný základ odlišný od jedné \textbf{prostá}. Znamená to, že různým hodnotám $x$
nemůže přiřadit stejnou funkční hodnotu.

Nejdůležitější vlastnost lze vyjádřit následovně:

\begin{definitionbox}
\[
    a^x = a^y \quad \Longleftrightarrow \quad x = y.
\]
\end{definitionbox}

Což lze říct jednoduše takto:  
\emph{Pokud se podaří exponenciální rovnici upravit do tvaru $a^{\text{výraz}_1} = a^{\text{výraz}_2}$, můžeme přejít 
k rovnici $\text{výraz}_1 = \text{výraz}_2$.}

\begin{examplebox}
\textbf{Příklad 1.}  
Upravme obě strany rovnice na stejný základ:
\[
    4^x = 16.
\]

Protože $4 = 2^2$ a $16 = 2^4$, přepíšeme rovnici ve tvaru
\[
    (2^2)^x = 2^4.
\]

Použitím pravidla pro mocninu mocniny získáme
\[
    2^{2x} = 2^4.
\]

Nyní mají oba výrazy stejný základ, a proto stačí porovnat exponenty:
\[
    2x = 4 \quad \Longleftrightarrow \quad x = 2.
\]
\end{examplebox}

\begin{notebox}
Tento postup ukazuje, že většina exponenciálních rovnic se snažíme upravit do tvaru
\[
    a^x = a^y,
\]
protože právě tehdy můžeme využít základní vlastnost exponenciální funkce:
\[
    a^x = a^y \;\Longleftrightarrow\; x = y.
\]
\end{notebox}

\begin{examplebox}
\textbf{Příklad 2.}  
Upravme rovnici tak, aby na obou stranách vystupoval stejný základ:
\[
    2 \cdot 2^x \cdot 8 = \frac{2^x \cdot 2^{x+1}}{8}.
\]

Nejprve vyjádříme všechny konstanty pomocí základu $2$.

Levá strana se tedy změní na
\[
    2^1 \cdot 2^x \cdot 2^3 = 2^{1 + x + 3} = 2^{x+4}.
\]

Pravou stranu rovnice upravíme obdobně:
\[
    \frac{2^x \cdot 2^{x+1}}{2^3}
    = \frac{2^{x + x + 1}}{2^3}
    = 2^{2x + 1 - 3}
    = 2^{2x - 2}.
\]

Tím dostáváme rovnici se stejným základem:
\[
    2^{x+4} = 2^{2x-2}.
\]

Protože se na obou stranách objevuje stejný základ, můžeme porovnat exponenty:
\[
    x + 4 = 2x - 2.
\]

Z čehož plyne:
\[
    x = 6.
\]

Tím jsme našli kořen (tedy $x$-ovou souřadnici průsečíku původních dvou výrazů):  
\[
K = \{6\}.
\]
\end{examplebox}

Důvod této ekvivalence je jednoduchý:

\begin{itemize}
    \item pro $a>1$ je funkce $a^x$ přísně \important{rostoucí},  
    \item pro $0<a<1$ je funkce $a^x$ přísně \important{klesající},
\end{itemize}

\noindent
a v obou případech tedy přiřazuje každé hodnotě $x$ právě jednu funkční hodnotu.

\subsubsection*{Důležitý fakt o oboru hodnot}

Exponenciální funkce má obor hodnot:
\[
H(f) = (0;\infty),
\]
což znamená, že $a^x$ je vždy kladné číslo.

Z toho plynou dvě zásadní skutečnosti:

\begin{itemize}
    \item výraz $a^x$ \textbf{nikdy} nemůže být roven nule ani zápornému číslu,
    \item rovnice $a^x = b$ má řešení \textbf{pouze tehdy, když $b>0$}.
\end{itemize}

Proto pokud $b \le 0$, rovnice $a^x = b$ \textbf{nemá řešení} a pokud $b > 0$, má rovnice \textbf{právě jedno řešení}.

Tento fakt tvoří základní východisko pro řešení rovnic s jedním exponenciálním výrazem.


\newpage
% ============================
%   M E T O D Y   Ř E Š E N Í
% ============================
\subsection{Metody řešení exponenciálních rovnic}

V některých exponenciálních rovnicích se vícekrát objevuje tentýž výraz,
například $2^x$. V takovém případě lze rovnici často řešit dvěma způsoby.

Je však důležité si uvědomit, že struktura rovnice zásadně ovlivňuje,
jaké úpravy jsou v dané situaci možné. Pokud se daný výraz objevuje
\textbf{pouze v násobení či dělení}, můžeme pomocí známých pravidel pro
mocniny jednoduše převádět výrazy na tvar
\[
    a^{f(x)} = a^{g(x)},
\]
protože součin $a^m \cdot a^n = a^{m+n}$ i podíl
$\dfrac{a^m}{a^n} = a^{m-n}$ nám dovolují přímo pracovat s exponenty.

Jakmile však v rovnici vystupuje \textbf{sčítání nebo odčítání exponenciálních výrazů},
situace se zásadně mění. Sčítání výrazů jako $2^x + 2^{x+1}$ nelze převést
na jedinou mocninu stejným způsobem, protože exponenty nelze jednoduše sečíst.
V takových rovnicích proto přímý přepis na tvar
$a^{f(x)} = a^{g(x)}$ není triviální ani možný.

Právě proto v těchto případech používáme dvě základní metody řešení:

\begin{itemize}
    \item \textbf{vytýkání společného exponenciálního výrazu},  
    \item \textbf{substituci}, která nahradí výraz typu $2^x$ novou proměnnou.
\end{itemize}

Obě metody mají stejný cíl — převést původní exponenciální rovnici
na jednodušší rovnici v jediné proměnné, kterou pak lze vyřešit
běžnými algebraickými postupy.


Obě metody vedou ke stejnému výsledku, ale každá je výhodná v jiné situaci.
Následující příklad ukazuje oba přístupy.


\begin{examplebox}
\textbf{Příklad 3.}  
Řešte rovnici:
\[
    3\cdot 2^x + 2^x = 32.
\]
\end{examplebox}

\noindent
Tato rovnice je ideální ukázkou, protože obsahuje dvakrát stejný výraz $2^x$.
Nejprve si ukážeme řešení pomocí vytýkání.

% ======================
% METODA 1: VYTÝKÁNÍ
% ======================

\begin{examplebox}
\textbf{Metoda 1: Řešení pomocí vytýkání}
 
Máme řešte rovnici:
\[
    3\cdot 2^x + 2^x = 32.
\]

Na levé straně se v obou sčítancích objevuje $2^x$, můžeme jej tedy vytknout:
\[
    3\cdot 2^x + 2^x = (3+1)\cdot 2^x = 4\cdot 2^x.
\]

Tím dostáváme jednodušší rovnici:
\[
    4\cdot 2^x = 32.
\]

Číslo $32$ vyjádříme jako mocninu základu $2$:
\[
    32 = 2^5.
\]

Také $4$ můžeme přepsat jako $2^2$:
\[
    2^2 \cdot 2^x = 2^{x+2} = 2^5.
\]

Nyní již mají obě strany stejný základ, můžeme proto porovnat exponenty:
\[
    x + 2 = 5 \quad \Longleftrightarrow \quad x = 3.
\]

Kořen rovnice je tedy $K=\{3\}$.
\end{examplebox}

\noindent
Metoda vytýkání je vhodná tehdy, když lze rychle rozpoznat společný exponenciální
činitel. V tomto případě vedla k velmi krátkému a přehlednému postupu.

\medskip

Nyní si ukážeme druhý způsob řešení – substituci.

% ======================
% METODA 2: SUBSTITUCE
% ======================

\begin{examplebox}
\textbf{Metoda 2: Řešení pomocí substituce}

Máme řešte rovnici:
\[
    3\cdot 2^x + 2^x = 32.
\]

Nejvhodnější substituce bude:
\[
    2^x = t, \qquad t>0.
\]

Tím převedeme původní rovnici na výrazně jednodušší lineární rovnici:
\[
    3t + t = 32 \quad \Longleftrightarrow \quad 4t = 32.
\]

Vyřešíme ji:
\[
    t = 8.
\]

Vrátíme substituci zpět:
\[
    2^x = 8
\]
\[
    2^x = 2^3
\]

Porovnáme exponenty:
\[
    x = 3.
\]

\medskip
\textbf{Výsledek:}  
Kořen rovnice je $K = \{3\}$.
\end{examplebox}

\noindent
Substituce je velmi univerzální metoda, která se hodí zejména v případech,
kdy se v rovnici vyskytuje více různých mocnin téhož základu (např. $2^x$, $2^{x+1}$, $4^x$ apod.).
V takových úlohách je substituce často přehlednější než vypisování všech mocnin
pomocí základního výrazu.




\begin{notebox}
\textbf{Kdy použít vytýkání?}  
Pokud se výraz $a^x$ vyskytuje ve více sčítancích a všechny je možné snadno upravit
na tvar $k \cdot a^x$, je vytýkání rychlejší a přehlednější metoda.

\medskip

\textbf{Kdy použít substituci?}  
Substituce je vhodná tehdy, když rovnice obsahuje více různých mocnin téhož výrazu,
například $2^x$, $2^{x+1}$, $4^x$, $8^{x-3}$ apod., a je obtížnější rychle odhadnout,
jak výraz vytknout. Substituce vše zjednoduší na rovnici v nové proměnné $t$.
\end{notebox}



\begin{examplebox}
\textbf{Příklad 4.}  
Řešte rovnici:
\[
    2^{x-1} \cdot \left(\frac{1}{3}\right)^{x-1}
    \;-\;
    2^x \cdot \left(\frac{1}{3}\right)^x
    \;=\;
    \frac{2}{9}.
\]

Na první pohled se zdá, že máme v rovnici hned dva exponenciální výrazy:
\[
a = 2^x 
\qquad\text{a}\qquad 
b = \left(\frac{1}{3}\right)^x.
\]

Taková substituce však \textbf{nedává smysl}.  
Pokud bychom zavedli dvě neznámé, objevily by se v rovnici dva parametry a dostali bychom
nekonečně mnoho dvojic $(a,b)$, které splňují původní vztah.  
Potřebujeme z obou výrazů vytvořit \textbf{jeden společný základ}.

\medskip
\textbf{1) Sjednotíme základy}

Využijeme toho, že
\[
    \left(\frac{1}{3}\right)^x = \frac{1^x}{3^x}
\]
A nyní využijeme poznatku z mocnin a sice, že $1^n = 1 $ a tedy můžeme ještě výraz upravit na:
\[
    \frac{1^x}{3^x} = \frac{1}{3^x}
\]
Obdobně pak urpavíme i
\[
    \left(\frac{1}{3}\right)^{x-1} = \frac{1^{x-1}}{3^{x-1}} = \frac{1}{3^{x-1}}
\]
To nyní dosadíme do rovnice:
\[
    2^{x-1} \cdot \frac{1}{3^{x-1}}
    \;-\;
    2^x \cdot \frac{1}{3^x}
    \;=\;
    \frac{2}{9}.
\]

To lze lépe přepsat jako:
\[
    \frac{2^{x-1}}{3^{x-1}}
    \;-\;
    \frac{2^x}{3^x}
    \;=\;
    \frac{2}{9}.
\]
Což ještě můžeme upravit na: 
\[
    \frac{2^{x}}{3^{x}} \cdot \frac{2^{-1}}{3^{-1}}
    \;-\;
    \frac{2^x}{3^x}
    \;=\;
    \frac{2}{9}.
\]

\[
    \frac{2^{x}}{3^{x}} \cdot \frac{3}{2}
    \;-\;
    \frac{2^x}{3^x}
    \;=\;
    \frac{2}{9}.
\]
\medskip
\textbf{2) Zavedeme substituci}

Nyní již můžeme bezpečně zavést substituci jediným výrazem:
\[
    t = \left(\frac{2}{3}\right)^x, \qquad t>0.
\]

Rovnice se změní na:
\[
    \frac{3}{2} t - t = \frac{2}{9}.
\]

Sečteme členy na levé straně:
\[
    \left(\frac{3}{2} - 1\right)t = \frac{2}{9}.
\]

Tedy:
\[
    \frac{1}{2}t = \frac{2}{9}
    \quad\Longrightarrow\quad
    t = \frac{4}{9}.
\]

\medskip
\textbf{3) Vrátíme substituci}

Máme:
\[
    \left(\frac{2}{3}\right)^x = \frac{4}{9}
    = \left(\frac{2}{3}\right)^2.
\]

Porovnáme exponenty:
\[
    x = 2.
\]

\medskip
\textbf{Výsledek:}  
Kořen rovnice je množina
\[
K = \{2\}.
\]
\end{examplebox}



\section{Soustavy exponenciálních rovnic}

Obdobně jako jsme v algebře pracovali se soustavami lineárních rovnic, kde jsme hledali takovou
hodnotu proměnné, která vyhovuje několika rovnicím současně, můžeme podobným způsobem postupovat
i u soustav rovnic exponenciálních. Zatímco v lineárních úlohách vystupovaly proměnné přímo
v základních tvarech výrazů, zde se objevují v exponentu, což s sebou přináší několik nových
postupů a technik. Základní idea však zůstává stejná: hledáme takovou hodnotu (nebo hodnoty)
proměnných, která splňuje všechny rovnice v soustavě zároveň.

\medskip

V exponenciálních soustavách se často setkáváme s výrazy jako $2^x$, $3^x$ nebo $4^y$, které
se v obou rovnicích vyskytují opakovaně. Abychom mohli soustavu účinně řešit, snažíme se buď
využít vlastností exponentů, případně sjednotit základy mocnin, nebo celou úlohu převést
substitucí na soustavu lineárních rovnic, kterou lze vyřešit známými metodami.

\medskip

Na rozdíl od soustav lineárních rovnic ovšem nemůžeme jednoduše sčítat nebo odčítat celé
rovnice, protože exponenciální výrazy se nedají spojovat přímo. Namísto toho využíváme
následující postupy:

\begin{itemize}
    \item vlastnost exponentů $a^x = a^y \Longleftrightarrow x = y$,  
    \item \textbf{vhodnou substituci}, která převede exponenciální výrazy na nové proměnné,
    \item úpravy rovnic pomocí sjednocení základů (např. přepis $4^y = 2^{2y}$),
    \item v některých úlohách i \textbf{dělení dvou rovnic}, pokud mají společný exponenciální výraz.
\end{itemize}

Nejčastěji se v soustavách exponenciálních rovnic objevují dvě situace:

\begin{enumerate}[label=\arabic*.]
    \item Obě rovnice obsahují stejný výraz (například $3^x$ nebo $4^y$), což umožní přímou substituci.
    \item Rovnice obsahují různé základy, které lze převést na společný základ a následně dosadit 
          jedinou substituci (například $8^x$ a $2^{3x}$).
\end{enumerate}

Následující příklad ukazuje typický postup, kdy správná volba substituce výrazně zjednoduší celou soustavu.

\begin{examplebox}
\textbf{Příklad 1.}  
Vyřešte soustavu rovnic:
\[
\begin{cases}
2 \cdot 3^x + 3 \cdot 4^y = 24, \\
3 \cdot 3^x - 10 \cdot 4^y = 7.
\end{cases}
\]

V obou rovnicích se objevují výrazy $3^x$ a $4^y$. Zavedením vhodné substituce tak 
můžeme převést soustavu na tvar velmi podobný běžné lineární soustavě.

\medskip
\textbf{Substituce:}  
\[
a = 3^x, \qquad b = 4^y.
\]

Po dosazení dostáváme soustavu lineárních rovnic, které již umíme řešit:
\[
\begin{cases}
2a + 3b = 24, \\
3a - 10b = 7.
\end{cases}
\]

\textbf{Úprava soustavy:}

Vynásobíme první rovnici číslem $3$:
\[
6a + 9b = 72.
\]

Druhou rovnici vynásobíme číslem $2$:
\[
6a - 20b = 14.
\]

Rovnice odečteme:
\[
(6a + 9b) - (6a - 20b) = 72 - 14,
\]
\[
29b = 58 \quad \Longrightarrow \quad b = 2.
\]

Dopočteme $a$:
\[
2a + 3b = 24,
\]
\[
2a + 6 = 24,
\]
\[
2a = 18 \quad \Longrightarrow \quad a = 9.
\]

\medskip
\textbf{Návrat k původní proměnné}

Máme
\[
a = 9 = 3^x.
\]
Protože $9 = 3^2$, dostáváme:
\[
3^2 = 3^x \quad\Longrightarrow\quad x = 2.
\]

Dále platí
\[
b = 2 = 4^y.
\]

Zapíšeme $4$ jako $2^2$:
\[
4^y = (2^2)^y = 2^{2y}.
\]

Z rovnice
\[
2^{2y} = 2
\]
plyne
\[
2y = 1 \quad\Longrightarrow\quad y = \frac{1}{2}.
\]

\medskip
Výsledek je tedy uspořádaná dvojice, která reperezentuje bod v kartézské soustavě prvků:
\[
K = \left\{\, \left[2;\ \frac{1}{2}\right] \right\}.
\]

\end{examplebox}



\newpage
\section{Exponenciální nerovnice}

Po vyřešení exponenciálních rovnic se přirozeně dostáváme k exponenciálním nerovnicím,
které řešíme velmi podobným způsobem. Základní rozdíl spočívá v tom, že zde nehledáme jedinou
hodnotu proměnné, ale takové hodnoty, pro něž je daný exponenciální výraz větší či menší než
zadané číslo nebo jiný exponenciální výraz. Přestože pracujeme s podobnými výrazy jako v rovnicích,
ujíždí se řešení nerovnic o jednu důležitou vlastnost exponenciální funkce – její monotónnost.

\medskip

Exponenciální funkce $y = a^x$ je pro každý základ $a>0$, $a\neq1$, prostá, a kromě toho je také
buď ryze rostoucí (pokud $a>1$), nebo ryze klesající (pokud $0<a<1$). Tato vlastnost přímo určuje,
jak se bude měnit směr nerovnosti při přechodu od tvaru
\[
a^{f(x)} \quad \text{vs.} \quad a^{g(x)}
\]
k porovnání samotných exponentů.

\begin{notebox}
\textbf{Monotónnost exponenciální funkce:}

\begin{itemize}
    \item Pokud $a > 1$, pak platí:  
    \[
        a^{f(x)} > a^{g(x)} \quad\Longleftrightarrow\quad f(x) > g(x).
    \]

    \item Pokud $0 < a < 1$, pak platí:  
    \[
        a^{f(x)} > a^{g(x)} \quad\Longleftrightarrow\quad f(x) < g(x),
    \]
    tedy se \textbf{směr nerovnosti obrací}.
\end{itemize}

\end{notebox}

Tato skutečnost je klíčová a tvoří základ veškerého řešení exponenciálních nerovnic.
Jakmile se podaří vyjádřit obě strany nerovnice stejným základem, můžeme místo exponenciální
nerovnice řešit obyčejnou nerovnici exponentů, a to se správným směrem podle uvedených pravidel.

\medskip

V praxi se nejčastěji setkáváme se třemi typy:

\begin{enumerate}[label=\arabic*.]
    \item Nerovnice se stejným základem:  
          \[
              a^{f(x)} > a^{g(x)}.
          \]

    \item Nerovnice s různými základy, které lze převést na jeden (např. $4 = 2^2$, $8 = 2^3$).

    \item Nerovnice tvaru $a^x > b$, kde $b>0$, které lze řešit přepisem $b$ na mocninu základu $a$.
\end{enumerate}

\begin{examplebox}
\textbf{Příklad 1.}  
Vyřešte nerovnici:
\[
    2^{3x-1} \; < \; 16.
\]

Nejprve převedeme pravou stranu na mocninu se stejným základem:
\[
16 = 2^4.
\]

Tím získáváme nerovnici
\[
    2^{3x-1} < 2^4.
\]

Protože základ $2>1$, exponenciální funkce je rostoucí, a proto můžeme porovnat
pouze exponenty \textit{bez změny směru nerovnosti} tedy:
\[
    3x - 1 < 4.
\]

Vyřešíme nerovnici:
\[
3x < 5 \quad\Longrightarrow\quad x < \frac{5}{3}.
\]

\medskip
\textbf{Řešení:}  
\[
K = \left(-\infty,\; \frac{5}{3}\right).
\]
\end{examplebox}

Nerovnici lze velmi dobře chápat i graficky: hledáme, pro které hodnoty proměnné leží graf exponenciální
funkce pod (nebo nad) hodnotou daného čísla, tedy kdy je křivka umístěna na jedné straně vodorovné přímky
$y = 16$.

\begin{center}
\begin{axisgraph}[
    xmin=-2, xmax=4,
    ymin=0, ymax=20,
    width=15.5cm,
    height=12cm,
    samples=300,
    domain=-2:4,
    restrict y to domain=0:20,
    unbounded coords=discard,
    title={Grafické znázornění nerovnice $2^{3x-1} < 16$},
]
    % Exponenciála
    \addplot[blue, thick] {2^(3*x - 1)};

    % Vodorovná přímka
    \addplot[red, dashed, very thick] {16};

    % Průsečík
    \addplot[
        mark=*,
        mark options={scale=1.4, fill=black},
    ] coordinates {(5/3,16)};

    % Svislá čára
    \addplot[dashed, gray] coordinates {(5/3,0) (5/3,16)};
\end{axisgraph}
\end{center}



\begin{notebox}
Výsledek nerovnice odpovídá těm hodnotám $x$, pro které leží graf exponenciály
$\;y = 2^{3x-1}\;$ pod vodorovnou přímkou $y=16$.  
Přesně v bodě $x = \frac{5}{3}$ se grafy dotýkají.
\end{notebox}






% -----------------------------------------------------------------
% ---------------------- P Ř Í K L A D Y --------------------------
% -----------------------------------------------------------------

\exercisepart
\setcounter{section}{1}

% =========================
% ÚLOHA 1 – Základní exponenciální rovnice
% =========================
\begin{exercise}[Základní exponenciální výrazy]
Vyřešte následující exponenciální výrazy
\begin{tasks}(3)
    \task $2^x = 4$
    \task $2^x = 8$
    \task $3^x = 9$
    \task $4^x = 16$
    \task $2^{2x} = 64$
    \task $3^x = \dfrac{1}{3}$
    \task $3^x = \frac{1}{81}$
    \task $5^x = 125$
    \task $5^x = 625$
    \task $3^x = \sqrt{3}$
    \task $7^{2x} = 49$
    \task $9^x = 3$
    \task $2^{x+1} = 16$
    \task $4^{x-2} = 1$
    \task $8^{x} = 2$
    \task $5^{x+1} = 25$
    \task $\left(\dfrac12\right)^{x} = 4$
    \task $2^{3x} = 32$
    \task $81^{-x} = 27$
    \task $4^{(x+3)(2-5x)} = 1$
    \task $3^{5-x} = 3^{4+x}$
\end{tasks}
\end{exercise}


% =========================
% ÚLOHA 2 – Exponenciální rovnice s úpravami mocnin
% =========================
\begin{exercise}[Exponenciální rovnice s úpravami mocnin]
Vyřešte následující exponenciální rovnice.

\begin{tasks}(3)
    \task $\dfrac{1}{5^{2x-4}}=125$
    \task $9^{x+1} \;=\; 3 \cdot 27^{x}$
    \task $\sqrt[x]{2\sqrt{8}} \;=\; \sqrt[x]{16}\cdot 2$
    \task $25^x = \left(\dfrac{1}{5}\right)^{x^2}$
    \task $\sqrt[3]{3^{2x}} \;=\; \dfrac{1}{3}$
    \task $\left(\dfrac{3}{4}\right)^{\frac{x}{3}} = \sqrt[3]{\left(\dfrac{4}{4}\right)^{x}}$
    \task $\dfrac{27}{8}\left(\dfrac{2}{3}\right)^{x} \;=\; \dfrac{4}{9}\left(\dfrac{9}{4}\right)^{x+2}$
    \task $\sqrt[\frac{1}{x}]{\dfrac{4^{4}}{16^{x}}} \;=\; 4\left(\dfrac{1}{2}\right)^{x}$
    \task $\left(\dfrac{1}{2}\right)^{x-1} \;=\; 8 \cdot \left(\dfrac{1}{2}\right)^{3x}$
    \task $4^{x-1} \cdot 8 \;=\; 2^{3x}$
    \task $5^{x+2} \;=\; 125 \cdot 25^{x}$
    \task $\sqrt{2^{x+1}} \;=\; 8$
    \task $7 \cdot 7^{x} \;=\; 343$
    \task $2 \cdot 2^x \cdot 4^{2-x} \;=\; \dfrac{8}{2^{3x+1}}$
    \task $\dfrac{9 \cdot 3^x}{\sqrt{3}\cdot \sqrt[3]{9^x}} \;=\; \dfrac{\sqrt[4]{3^x}}{27}$
    \task $2 \cdot 2^x \cdot 8 \;=\; \dfrac{2^x \cdot 2^{x+1}}{8}$
    \task $\dfrac{3^x}{9^{x-2}} \;=\; \dfrac{27^{x}}{9 \cdot 3^{4-x}}$
    \task $2^x \cdot 5^x = 0,1 \cdot \left(10^{x-1}\right)^5$
    \task $\star \; \dfrac{6^{x^2}}{2^{-15}} = \dfrac{3^{-15}}{6^{12-12x}}$
    \task \text{$\star \; 0,25^{2-\sqrt{5x+1}} = 4 \cdot 2^{\sqrt{5x+1}}$}
    \task $\star \; 9^{x(x-1) - 0,5} = \sqrt{3}$
\end{tasks}
\end{exercise}


% =========================
% ÚLOHA 3 – Exponenciální rovnice (pokročilejší úpravy)
% =========================
\begin{exercise}[Pokročilé exponenciální rovnice]
Vyřešte následující rovnice.

\begin{tasks}(2)
\task $2^{3x-1} \cdot 4 = 8^{x+1} \cdot \left(\dfrac{1}{2}\right)^{x}$
\task $\sqrt[4]{4^{x}} \cdot \sqrt[3]{2^{x-3}} = \sqrt[6]{16}$
\task $\dfrac{1}{3^{x}}
      = \dfrac{1}{\sqrt{3}} \cdot \sqrt[6]{27^{3 - 3x}} \cdot \left(\dfrac{1}{9}\right)^{x+3}$
\task $0.25 \cdot \left(\dfrac{1}{4}\right)^{2x} = 1$
\task $2 \cdot 0.5^{\,x + \tfrac{8}{3}x} = \dfrac{8}{\sqrt[3]{4}}$
\task $5^{x} \cdot 2^{x} = 100^{x-1}$
\task $\dfrac{81}{16}
      = \left(\dfrac{2}{3}\right)^{x} \cdot \left(\dfrac{9}{4}\right)^{x+1}$
\task $\dfrac{3^{x}}{2 \cdot 3^{\sqrt{3}}} = 4.5$

\end{tasks}

\end{exercise}



% =========================
% ÚLOHA 4 – Exponenciální rovnice s kombinovanými výrazy
% =========================
\begin{exercise}[Kombinované exponenciální rovnice]
Vyřešte následující rovnice pomocí vytýkání.

\begin{tasks}(2)
    \task $3 \cdot 2^{x} + 2^{x} = 32$                    %  x = 3
    \task $2 \cdot 3^{x} + 3^{x+1} = 5$                   %  x = 0
    \task $\dfrac{3^{x+4}}{9} + 9 \cdot 3^{x} = \dfrac{2}{3}$   %  x = -3
    \task $3^{x} - 2 \cdot 3^{x-1} + 3^{x-2} - 2 \cdot 3^{x-3} = 30$  % x = 4
    \task $2^{2-x} = 2^{4-x} - 3\sqrt[3]{2}$              %  x = \tfrac{5}{3}
    \task $2 \cdot 4^{-x} - 9 \cdot 2^{x} + 4 = 0$        %  (tu klidně můžeš nechat jako těžší/na diskusi, ale řešení je reálné)
    \task $4^{x} + 4^{\frac{3}{2}-x} = 9$                %  x = 0 nebo x = \tfrac{3}{2}
    \task $5 \cdot 2^{x+1} - 3 \cdot 2^{x} = 56$          %  x = 3
    \task $3^{2x} - 3^{x+1} = 54$                         %  x = 2
    \task $7^{x} + 7^{x-1} = 56$                          %  x = 2
    \task $9^{x} - 3^{x+1} = 54$                          %  x = 2
    \task $2^{x} + 4^{x} = 20$                            %  x = 2
    \task $3^{x} + 3^{-x} = \dfrac{10}{3}$                %  x = 1 nebo x = -1
    \task $5^{x} - 5^{2-x} = 0$                           %  x = 1
    \task $2^{x+2} + 2^{1-x} = 6$                         %  x = 0 nebo x = -1
    \task $2^{v} + 2^{v+1} = 24$
    \task $9^{v+2} + 5\cdot9^{v+1} = 14$
    \task $3^{x} + 3^{x+1} = 108$
    \task $3^{x+2} + 9^{x+1} = 810$
    \task $2^{x+1} + 2^{x-1} + 2^{x+3} = \dfrac{21}{8}$
    \task $7 \cdot 4^{-x+2} = 3 \cdot 4^{-x+3} - 5$
    \task $\dfrac{4}{5} \cdot 5^{0} + 5^{-1} - 25^{x} + 20 \cdot 25^{x-1} = 0$
    \task $3^{x} \cdot \left(\dfrac{1}{2}\right)^{x} + 3^{x+1} \cdot \left(\dfrac{1}{2}\right)^{x+1} = \dfrac{5}{3}$
    \task $2^{2x} \cdot 5^{x} - 2^{2x-1} \cdot 5^{x+1} = -600$
\end{tasks}

\end{exercise}



% =========================
% ÚLOHA 6 – Exponenciální rovnice řešené substitucí
% =========================
\begin{exercise}[Kombinované exponenciální rovnice]
Vyřešte následující exponenciální rovnice pomocí substituce.

\begin{tasks}(2)
    \task $2^{x-1}\left(\dfrac{1}{3}\right)^{x-1} - 2^{x}\left(\dfrac{1}{3}\right)^{x} = \dfrac{2}{9}$
    \task $20 \cdot 2^{x} - 2^{x+1} = 3^{x+2} - 3^{x}$
    \task $3 \cdot 5^{x-2} + 7 \cdot 5^{x-3} = 5 \cdot 5^{x-3} + 5^{x-1}$
    \task $3 \cdot 2^{2-x} + 2 \cdot 2^{1-x} = 8 \cdot 3^{2-x} + 3 \cdot 3^{1-x}$
    \task $2 \cdot 4^{x} - 3^{x-\tfrac12} = 3^{x+\tfrac12} + 2 \cdot 2^{2x-1}$
    \task $9^{x+1} + 5 \cdot 6^{x} = 4^{x+1}$
    \task $4 \cdot 16^{x} + 4^{x+1} = 2 \cdot 8^{x+1} + 2^{3x}$
    \task $2^{4x} - 50 \cdot 2^{2x} = 896$
    \task $4^x + 2^x - 6 = 0$
    \task $3^{x+2} + 9^{x+1} = 810$
    \task $3^x + 3^{x+1} - 5^{x+1} = 5^x - 3^{x+3} + 5^{x+2}$
    \task $\dfrac{2^{x+3} \cdot 3^{x+2}}{6^{7-x} \cdot 8^{x-1}} = \dfrac{9^{x-2}}{3}$
\end{tasks}

\end{exercise}


\begin{exercise}[Exponenciální rovnice]
Vyřešte následující rovnice pomocí substituce.
\begin{tasks}(2)

    \task $4^{2x} - 6 \cdot 4^{x} + 8 = 0$
    \task \(\dfrac{1}{4} \cdot 2^{x} + \dfrac{1}{2} \cdot 4^{x} = 9\)
    \task $9 \cdot 3^{x} + 3^{-x} = 10$
    \task $3^{2x} - 4 \cdot 3^{x} = -3$
    \task $2^{x+3} + 2^{3-x} = 20$
    \task $4^{x} + 2^{x+1} = 24$
    \task $5^{x-1} + 5^{1-x} = 2$
    \task $7^{2x} - 14 \cdot 7^{x} + 49 = 0$
    \task $9^{x+1} - 27 \cdot 3^{x} = 0$    
    \task $9^{x-0{,}5} + 9^{0{,}5-x} = \dfrac{10}{3}$
    \task $2\left(\dfrac{1}{4}\right)^{x}
           - 3\left(\dfrac{1}{2}\right)^{x}
           = \bigl[\,1 + \left(\dfrac{1}{2}\right)^{x}\bigr]
             \left(\dfrac{1}{4}\right)^{-1}$
    \task $5^{2x}\bigl(5^{2x} - 5\bigr)
           = 3\bigl(5^{2x+1} + 5^{2x}\bigr) + 50$
\end{tasks}
\end{exercise}


\setcounter{section}{2}
% =========================
% ÚLOHA 8 – Základní exponenciální nerovnice
% =========================
\begin{exercise}[Základní exponenciální nerovnice]
Řešte následující nerovnice s neznámou $x \in \mathbb{R}$.

\begin{tasks}(3)
    \task $3^{x+5} < 1$
    \task $2^{3x-4} \ge 1$
    \task  $0{,}1^{\,2x} \le 1$
    \task $\left(\dfrac{1}{2}\right)^{x+3} > 1$
    \task $3^{x-5} < 0$
    \task $\left(\dfrac{1}{5}\right)^{3x+2} \ge 0$
\end{tasks}

\end{exercise}


% =========================
% ÚLOHA 9 – Exponenciální nerovnice s různými základy
% =========================
\begin{exercise}[Exponenciální nerovnice s různými základy]
Řešte následující nerovnice s neznámou $x \in \mathbb{R}$.

\begin{tasks}(3)
    \task $3^{x-1} > 9$
    \task $0{,}125^{\,2x+3} \le 16$
    \task $3^{2x} \le 7$
    \task $0{,}2^{\,x} > 3$
    \task $4^{x} \cdot 2^{x} \le 100$
    \task $\left(\dfrac{2}{7}\right)^{x} \cdot 3^{x+4} > -1$
\end{tasks}

\end{exercise}


% =========================
% ÚLOHA 10 – Porovnání exponentů v nerovnicích
% =========================
\begin{exercise}[Exponenciální nerovnice – porovnání exponentů]
Řešte následující nerovnice s neznámou $x \in \mathbb{R}$.

\begin{tasks}(3)
    \task $3^{x+4} \ge 3^{2x-1}$
    \task $4^{x+5} < 16^{x+1}$
    \task $0{,}2^{\,3x-2} \ge 0{,}2^{\,x+5}$
    \task $\left(\dfrac{1}{6}\right)^{x^{2}} < \left(\dfrac{1}{6}\right)^{2x+8}$
    \task $2^{x} \cdot 2^{x+1} < 2^{x+3}$
    \task $\left(\dfrac{1}{2}\right)^{x} \cdot \dfrac{1}{2} < \dfrac{1}{8}$
\end{tasks}
\end{exercise}


\setcounter{section}{3}
% =========================
% ÚLOHA 11 – Soustavy exponenciálních rovnic
% =========================
\begin{exercise}[Soustavy exponenciálních rovnic]
Vyřešte následující soustavy rovnic. Vhodně použijte substituci.

\begin{tasks}(2)
    \task 
    $\begin{cases}
        2^{x} + 5 \cdot 3^{y} = 53,\\[2mm]
        7 \cdot 2^{x} - 3^{y} = 47
    \end{cases}$
    \task
    $\begin{cases}
        2^{x+1} + 3^{y} = 31,\\[2mm]
        2^{x} - 3^{y-2} = -1
    \end{cases}$
    \task
    $\begin{cases}
        2 \cdot 2^{x-y} + 2^{x+y-1} = 20,\\[2mm]
        10 \cdot 2^{x-y-1} - 2^{x+y} = -22
    \end{cases}$
    \task
    $\begin{cases}
        16^{x+y} = 8,\\[2mm]
        16^{x} + 16^{y} = 6
    \end{cases}$
    \task
    $\begin{cases}
        5^{\frac{x+y}{2}} : 5^{\frac{y}{4}} = 625,\\[2mm]
        6^{\frac{x-y}{2}} : 6^{\frac{x}{4}} = 216
    \end{cases}$
    \task
    $\begin{cases}
        x^{y+1} = 27,\\[2mm]
        x^{-y+1} = 3^{-1}
    \end{cases}$
    \task
    $\begin{cases}
        2 \cdot 3^{x} + 3 \cdot 4^{y} = 24,\\[2mm]
        3 \cdot 3^{x} - 10 \cdot 4^{y} = 7
    \end{cases}$
    \task
    $\begin{cases}
        0{,}5^{\,x} + 3^{y-1} = 25,\\[2mm]
        2 \cdot 0{,}5^{x+3} - 3^{y-2} = 1
    \end{cases}$
    \task
    $\begin{cases}
        81^{x+y} = 27,\\[2mm]
        81^{x} + 81^{y} = 12
    \end{cases}$
    \task
    $\begin{cases}
        2^{x+y} + 2^{x-y-1} = 3^{2},\\[2mm]
        2^{x+y+1} - 3 \cdot 2^{x-y+1} = 2^{2}
    \end{cases}$
\end{tasks}

\end{exercise}


\setcounter{section}{4}
% =========================
% ÚLOHA 12 – Parametry exponenciální funkce
% =========================
\begin{exercise}[Určení parametrů z daných bodů]
Určete reálná čísla $a,b$ tak, aby graf dané exponenciální funkce procházel
zadanými body v rovině.

\begin{tasks}(2)
    \task Určete reálná čísla $a,b$ tak, aby graf funkce
    \[
        y = 2^{x+a} + b
    \]
    procházel body $[-1;-2]$ a $[2;12]$.
    \task Určete reálná čísla $a,b$ tak, aby graf funkce
    \[
        y = 3^{x+a} + b
    \]
    procházel body $[0;5]$ a $[2;13]$.
\end{tasks}

\end{exercise}


% =========================
% ÚLOHA 13 – Shrnutí: mix exponenciálních rovnic
% =========================
\begin{exercise}[Shrnující exponenciální rovnice]
Vyřešte následující rovnice pomocí vhodných postupů.

\begin{tasks}(2)
    \task $\left(\dfrac{8}{27}\right)^{x} = \dfrac{9}{4}$
    \task $4 \cdot 2^{x} \cdot \sqrt{2} = 4^{x} \cdot 2$
    \task $3^{x+1} + 2 \cdot 3^{x} = 5^{x+1} - 2 \cdot 5^{x}$
    \task $3^{x} - 3^{x-1} - 2 \cdot 3^{x-2} - 3 \cdot 3^{x-3} = 9$
    \task $\dfrac{5^{x^{2}} \cdot 2^{x^{2}}}{10^{x}} = 1$
    \task $2 \cdot 4^{x} + 3 \cdot 9^{x} = 5 \cdot 6^{x}$
    \task $\sqrt[x]{2^{3}} \cdot 8 = \sqrt[x]{2} \cdot 4$
    \task $2^{2x} \cdot 5^{x+1} = 4^{x+1} + 4^{x}$
    \task $3^{x} + 2 \cdot 3^{1-x} = 5$
    \task
        $\begin{cases}
            2^{x+1} + 3 \cdot 3^{y} = 42,\\[2mm] 
            4 \cdot 2^{x} - 9^{\,y-1} = 4.
        \end{cases}$
\end{tasks}

\end{exercise}



\end{document}
